\chapter{Technische Erprobung: Proof of Concept}
\label{chap:poc}

Bevor eine Plattform ausgewählt wird, werden die beiden zentralen quelloffenen Kandidaten praktisch
erprobt. Ziel der technischen Erprobung ist es, die grundsätzliche Machbarkeit eines selbst
betriebenen Aufbaus zu überprüfen. Im Fokus steht kein Produktivbetrieb, sondern ein
nachvollziehbarer Proof of Concept (PoC), der zeigt, ob und wie sich die Plattformen auf einem Host
deployen und betreiben lassen. Die Ergebnisse fließen anschließend in die Bewertung in
Kapitel~\ref{chap:auswahl} ein.

\section{Proof of Concept: CyberRangeCZ}
\label{sec:poc-crcz}

Der PoC für CyberRangeCZ demonstriert zentrale Abläufe der Plattform: Installation und Deployment,
Erstzugriff über die Weboberflächen, Import einer Trainingsdefinition, Bereitstellung einer Sandbox
sowie den Zugriff auf die bereitgestellten Systeme. Verwendet wurde die CyberRangeCZ-Lite-Variante,
da sie einen schnellen Einstieg in die Plattformmechanik ermöglicht und zugleich die wesentlichen
Konzepte (Portal, Trainingsdefinitionen, Pools, Sandboxes) abbildet \cite{CyberRangeCZDevopsLite}.

\subsection{Testumgebung}
\label{sec:poc-testumgebung}

Der PoC wurde auf einer dedizierten Maschine mit Ubuntu 24.04 LTS durchgeführt. Die verwendeten
Ressourcen (CPU, RAM und Storage) sind in Abbildung~\ref{fig:poc-hardware} dargestellt.

\begin{figure}[H]
    \centering
    \includegraphics[width=0.85\linewidth]{poc_hardware.png}
    \caption{Testumgebung für den PoC (eigene Aufnahme)}
    \label{fig:poc-hardware}
\end{figure}

\subsection{Installation und Deployment der PoC-Umgebung}
\label{sec:poc-installation}

Die Installation folgt dem Ansatz, die notwendigen Virtualisierungs- und Container-Komponenten lokal
bereitzustellen und das eigentliche Deployment über das bereitgestellte PoC-Repository automatisiert
auszuführen. Dazu wurden zunächst Basiswerkzeuge sowie KVM/libvirt und Docker installiert.
KVM/libvirt stellt die Virtualisierungsbasis bereit (VM-Lifecycle und Netzwerk), während Docker
genutzt wird, um eine vorbereitete Umgebung auszuführen, in der Vagrant mit libvirt-Anbindung zur
Provisionierung verwendet wird.

\begin{verbatim}
sudo apt install -y \
  git tmux curl \
  qemu-kvm libvirt-daemon-system libvirt-clients bridge-utils cpu-checker \
  docker.io
\end{verbatim}

Im nächsten Schritt wurden die relevanten Dienste aktiviert und der Benutzer für die Verwendung von
\texttt{libvirt}, KVM und Docker berechtigt. Damit sind Virtualisierung und Containerbetrieb ohne
ständige Root-Kontexte nutzbar.

\begin{verbatim}
sudo systemctl enable --now libvirtd docker
sudo usermod -aG libvirt,kvm,docker $USER
newgrp docker
\end{verbatim}

Anschließend wurde das PoC-Repository geklont und die libvirt-URI gesetzt, sodass Vagrant eindeutig
die System-libvirt-Instanz adressiert. Wegen der langen Laufzeit und der vielen Ausgaben wurde der
Prozess in einer tmux-Session gestartet.

\begin{verbatim}
git clone https://github.com/cyberrangecz/devops-crczp-lite.git
cd devops-crczp-lite
export LIBVIRT_DEFAULT_URI=qemu:///system
tmux new -s crczp
\end{verbatim}

Das eigentliche Deployment erfolgt über einen Docker-Container, der Vagrant mit libvirt-Unterstützung
bereitstellt. Die relevanten Sockets und Verzeichnisse werden in den Container gemountet, damit
Vagrant (1) auf libvirt zugreifen kann, (2) seine Vagrant-States dauerhaft speichern kann und (3) im
aktuellen Arbeitsverzeichnis die Definitionsdateien findet. Der Parameter \texttt{--network host}
stellt sicher, dass Netzwerkoperationen und lokale Zugriffe (z.\,B.\ auf die Weboberfläche) ohne
zusätzliche Container-NAT-Komplexität funktionieren.

\begin{verbatim}
docker run -it --rm \
  -e LIBVIRT_DEFAULT_URI \
  -v /var/run/libvirt/:/var/run/libvirt/ \
  -v ~/.vagrant.d:/.vagrant.d \
  -v $(realpath "${PWD}"):${PWD} \
  -w "${PWD}" \
  --network host \
  vagrantlibvirt/vagrant-libvirt:latest \
  vagrant up
\end{verbatim}

Nach erfolgreichem Abschluss wurden eine URL für das Web-Portal sowie initiale Zugangsdaten
ausgegeben. Diese Ausgabe ist in Abbildung~\ref{fig:poc-deploy} dokumentiert.

\begin{figure}[H]
    \centering
    \includegraphics[width=0.92\linewidth]{poc_deploy_done.png}
    \caption{Erfolgreiches Deployment der PoC-Umgebung (eigene Aufnahme)}
    \label{fig:poc-deploy}
\end{figure}

\subsection{Erstzugriff über Keycloak und CyberRangeCZ-Portal}
\label{sec:poc-webzugriff}

CyberRangeCZ verwendet für Authentifizierung und Benutzer- und Client-Verwaltung typischerweise
Keycloak als Identity- und Access-Management-Komponente. Der erfolgreiche Zugriff auf die
Keycloak-Verwaltungskonsole bestätigt, dass die Authentifizierungs- und Basisdienste funktionsfähig
bereitgestellt wurden (Abbildung~\ref{fig:poc-keycloak}).

\begin{figure}[H]
    \centering
    \includegraphics[width=0.95\linewidth]{poc_keycloak.png}
    \caption{Keycloak-Verwaltungskonsole der PoC-Installation (eigene Aufnahme)}
    \label{fig:poc-keycloak}
\end{figure}

Der zentrale Einstiegspunkt für Trainings- und Sandbox-Management ist das CyberRangeCZ Platform
Portal. Abbildung~\ref{fig:poc-portal} zeigt die Startseite mit den wichtigsten Funktionen, unter
anderem Trainingsdefinitionen, Trainingsinstanzen, Sandboxes, Pools sowie Benutzer- und
Gruppenverwaltung.

\begin{figure}[H]
    \centering
    \includegraphics[width=0.95\linewidth]{poc_portal_home.png}
    \caption{CyberRangeCZ Platform Portal (Startseite) nach erfolgreichem Login (eigene Aufnahme)}
    \label{fig:poc-portal}
\end{figure}

\subsection{Import einer Trainingsdefinition und Provisionierung einer Sandbox}
\label{sec:poc-training-import}

Um die Kernfunktionalität der Plattform zu validieren, wurde eine Beispiel-Trainingsbibliothek
(\texttt{library-demo-training}) aus einem Git-Repository in das Portal eingebunden
\cite{CyberRangeCZLibraryDemoTraining}. Nach dem Import wurde ein Pool beziehungsweise eine
Trainingsinstanz erstellt und die Sandbox-Provisionierung ausgelöst. Die resultierende
Netzwerktopologie der bereitgestellten Umgebung ist in Abbildung~\ref{fig:poc-topology} dargestellt.

\begin{figure}[H]
    \centering
    \includegraphics[width=0.80\linewidth]{poc_topology.png}
    \caption{Netzwerktopologie einer bereitgestellten Demo-Sandbox (eigene Aufnahme)}
    \label{fig:poc-topology}
\end{figure}

\subsection{Zugriff auf die Sandbox und funktionaler Verbindungstest}
\label{sec:poc-zugriff}

Für den Zugriff auf die bereitgestellten Systeme stellt CyberRangeCZ im PoC automatisch erzeugte
SSH-Artefakte zur Verfügung. Dazu gehören typischerweise eine Benutzer-Konfiguration sowie ein
passender privater und öffentlicher Schlüssel. Abbildung~\ref{fig:poc-ssh-files} zeigt die
entsprechenden Dateien, die nach der Sandbox-Provisionierung bereitgestellt wurden.

Neben SSH sind in Cyber-Range-Umgebungen grundsätzlich auch weitere Zugriffsmöglichkeiten relevant,
insbesondere wenn Trainings- oder Forschungsszenarien eine grafische Interaktion erfordern. Je nach
Szenariodefinition kann der Zugriff daher auch über browserbasierte Oberflächen oder GUI-Zugriffe
erfolgen, etwa über eine Web-Konsole oder Remote-Desktop-Mechanismen wie RDP, VNC oder noVNC. Welche
Zugriffspfade tatsächlich verfügbar sind, hängt von der jeweiligen Trainingsdefinition und der
bereitgestellten Infrastruktur ab.

\begin{figure}[H]
    \centering
    \includegraphics[width=0.90\linewidth]{poc_ssh_files.png}
    \caption{SSH-Daten für den Sandbox-Zugriff (eigene Aufnahme)}
    \label{fig:poc-ssh-files}
\end{figure}

Zur Verifikation, dass der Zugriff auf die Sandbox tatsächlich funktioniert und Netzwerkkommunikation
möglich ist, wurde aus der bereitgestellten Client-Umgebung ein einfacher Scan gegen das
Server-System durchgeführt. Abbildung~\ref{fig:poc-nmap} zeigt exemplarisch, dass das Zielsystem
erreichbar und SSH (Port 22) offen ist. Damit ist die grundlegende Nutzbarkeit der bereitgestellten
Umgebung für weitere Schritte nachgewiesen.

\begin{figure}[H]
    \centering
    \includegraphics[width=0.92\linewidth]{poc_nmap.png}
    \caption{Einfacher Funktionstest innerhalb der Sandbox (eigene Aufnahme)}
    \label{fig:poc-nmap}
\end{figure}

\subsection{Erkenntnisse aus dem PoC}
\label{sec:poc-erkenntnisse}

Der PoC bestätigt, dass CyberRangeCZ in der Lite-Variante in einer lokalen Umgebung installiert und
betrieben werden kann und dabei zentrale Plattformmechanismen sichtbar werden: versionierbare
Trainingsdefinitionen, Portal-gestützte Verwaltung sowie automatisierte Provisionierung isolierter
Umgebungen. Besonders relevant ist, dass sich eine Trainingsdefinition aus einem Git-Repository
importieren und anschließend reproduzierbar als Sandbox instanziieren lässt.

Für die weitere Arbeit ergibt sich daraus als nächster Schritt, die Erprobung in Richtung
Tool-Integration zu vertiefen. Konkret bedeutet das, Mess- und Beobachtungspunkte (Logs, Metriken,
Netzwerkdaten) in einem Use Case gezielt zu platzieren und den Aufwand für die Integration und den
Vergleich mehrerer Sicherheits- oder Monitoring-Werkzeuge realistisch abzuschätzen.

\section{Proof of Concept: OCR}
\label{sec:poc-ocr}

Neben CyberRangeCZ wurde mit Open Cyber Range (OCR) eine zweite quelloffene Plattform praktisch
evaluiert. Im Vordergrund stand auch hier kein Produktivbetrieb, sondern ein nachvollziehbarer PoC,
der zeigt, ob sich OCR auf einem Host deployen lässt und wie die Architektur in der Praxis handhabbar
ist.

Im Unterschied zu vielen ``All-in-one''-Trainingsplattformen ist OCR nicht als monolithisches
Webprodukt konzipiert, sondern als Baukastensystem aus mehreren Diensten, die gemeinsam den
Lebenszyklus einer Übung abdecken: (1) Verwaltung und Distribution von Übungsartefakten, (2)
Orchestrierung und Zustandsverwaltung von Deployments, (3) Anbindung an eine konkrete Infrastruktur
(Virtualisierung oder Cloud) sowie (4) Identity- und Access-Management. Diese modulare Architektur
ist für die Forschung prinzipiell attraktiv, erhöht im Self-Hosted-Betrieb jedoch den
Integrationsaufwand erheblich, da mehrere Komponenten konsistent konfiguriert und miteinander
verbunden werden müssen \cite{OCRDeputyRepo,OCRCZRangerDoc}.

Für den PoC wurden die Kernkomponenten Deputy und Ranger betrachtet. Deputy fungiert als digitale
Bibliothek für Szenarioinhalte. Artefakte wie Packages, Templates oder vorbereitete
Übungskomponenten werden dort versioniert abgelegt und können von anderen OCR-Komponenten abgerufen
werden \cite{OCRDeputyRepo,OCRCZFirstDeputyPackage}.

Die Ranger-Komponente bildet die Control Plane für die Ausführung. Sie verwaltet den Zustand einer
Übung und koordiniert das Ausrollen in eine Zielumgebung. Entscheidend ist dabei, dass Ranger die
Infrastruktur nicht selbst erzeugt, sondern die Provisionierung an umgebungsspezifische Handler und
Deployer delegiert. Für OpenStack existiert dafür ein eigener Handler als separates Projekt
\cite{OCROpenstackHandler}.

\subsection{Testumgebung}
\label{sec:poc-ocr-testumgebung}

Der PoC wurde auf einer dedizierten Maschine mit Ubuntu 24.04 LTS durchgeführt. Die verwendeten
Ressourcen (CPU, RAM und Storage) sind in Abbildung~\ref{fig:poc-ocr-hardware} dargestellt.

\begin{figure}[H]
    \centering
    \includegraphics[width=0.85\linewidth]{poc_ocr_hardware.png}
    \caption{Testumgebung für den OCR-PoC (eigene Aufnahme)}
    \label{fig:poc-ocr-hardware}
\end{figure}

\subsection{Installation und Deployment der PoC-Umgebung}
\label{sec:poc-ocr-installation}

Die Installation orientierte sich an der offiziellen Cyber-Range-Dokumentation für die
OCR-Komponenten Deputy und Ranger \cite{OCRCZRangerDoc,OCRCZFirstDeputyPackage}. In der praktischen
Umsetzung traten allerdings viele Probleme auf, die nicht in der Dokumentation standen. Der
Ranger-Service prüft beim Start die Konfiguration zum Beispiel strikt auf Vollständigkeit und
korrekte Datentypen. Fehlende Pflichtfelder, etwa für Deployer- und Deployment-Strukturen, führten
daher zu wiederholten Startabbrüchen und Neustartschleifen.

\subsubsection{Deployment von Deputy}
\label{subsec:poc-ocr-deputy}

Deputy wurde als erster Baustein gestartet, da Ranger für den späteren Betrieb typischerweise
Artefakte über Deputy bezieht. Der Start erfolgte über Docker Compose.
Abbildung~\ref{fig:poc-ocr-deputy-up} dokumentiert den erfolgreichen Image-Pull und die laufenden
Container.

\begin{verbatim}
cd /opt/ocr/deputy
docker compose up -d
docker ps
\end{verbatim}

\begin{figure}[H]
    \centering
    \includegraphics[width=1\linewidth]{poc_ocr_deputy_up.png}
    \caption{Deputy: erfolgreicher Start der Services (eigene Aufnahme)}
    \label{fig:poc-ocr-deputy-up}
\end{figure}

\subsubsection{Deployment der Ranger-Komponente}
\label{subsec:poc-ocr-ranger}

Im nächsten Schritt wurde die Ranger-Komponente deployt. In der Praxis erwies sich dieser Teil als
deutlich aufwändiger als erwartet, da Ranger bereits im Minimalbetrieb strikte Anforderungen an die
Konfiguration stellt. Ranger validiert beim Start eine Reihe verpflichtender Felder und Datentypen.
Fehlen diese Angaben oder passen sie nicht exakt, beendet sich der Dienst sofort mit Fehlermeldungen.

Der Start und der Containerstatus sind in Abbildung~\ref{fig:poc-ocr-ranger-up} dokumentiert.
Abbildung~\ref{fig:poc-ocr-ranger-logs} zeigt den erfolgreichen Start über die Logmeldungen des
Webservers.

\begin{figure}[H]
    \centering
    \includegraphics[width=1\linewidth]{poc_ocr_ranger_up.png}
    \caption{Ranger-Komponente (eigene Aufnahme)}
    \label{fig:poc-ocr-ranger-up}
\end{figure}

\begin{figure}[H]
    \centering
    \includegraphics[width=1\linewidth]{poc_ocr_ranger_logs.png}
    \caption{Ranger-Komponente: erfolgreicher Start (eigene Aufnahme)}
    \label{fig:poc-ocr-ranger-logs}
\end{figure}

Im Rahmen des Vorprojekts konnte der PoC für OCR in der verfügbaren Zeit nicht bis zu einem
lauffähigen Trainings- oder Deployment-Start abgeschlossen werden. Zentrale Komponenten (Deputy und
Ranger) konnten zwar grundsätzlich gestartet werden, doch die praktische Integration der modularen
OCR-Architektur erwies sich als deutlich komplexer als erwartet, da mehrere Dienste, unter anderem
das Identity- und Access-Management, konfiguriert und miteinander abgestimmt werden müssen.
Zusätzlich setzt die eigentliche Bereitstellung von Übungsumgebungen eine angebundene
Virtualisierungs- oder Cloud-Infrastruktur voraus, etwa eine eigene OpenStack-Umgebung oder
alternativ VMware/ESXi über passende Handler und Deployer. Der PoC bleibt daher auf den erfolgreichen
Aufbau der Basisdienste und die technische Erkundung der Architektur beschränkt.